% Evans second edition, section 5.5.
% English source: printed p. 256; PDF page 273 (1-based).
% Chinese section starts checked against the chapter opening on PDF page 187; see docs/CHAPTER-05.md.
% Status: editorial study draft; expanded proofs and Chinese mathematical annotations.
\section{Traces}
\label{sec:5-5}

\subsection{Boundary values as a continuous operator}
\begin{theorem}[Evans \S5.5, Theorem 1]\label{thm:5-trace}
If $U$ is bounded with $C^1$ boundary and $1\leq p<\infty$, there is a
unique bounded linear map
\[
 T:W^{1,p}(U)\longrightarrow L^p(\partial U)
\]
agreeing with restriction for functions in $W^{1,p}(U)\cap C(\overline U)$.
In particular, $\|Tu\|_{L^p(\partial U)}\leq C\|u\|_{W^{1,p}(U)}$.
\end{theorem}
\Needspace{45mm}
\PDEAnnotation{边界不能随意赋值}{改变 $u$ 在零测集上的值不改变 Sobolev 元素；边界通常正是体积零测集。迹由域内逼近的极限决定，不能通过任意指定边界代表值来定义。}
\begin{expandedproof}

\PDEMarginDiagram{5.M}{沿法向控制边界值}
{先沿固定 $x'$ 的线段比较内部值与边界值，再对高度和切向变量积分。图中外法向朝下。}
{ch05-trace-strip-margin}

\paragraph{1. A flat strip estimate.}
For smooth $v$ on a strip $Q'\times[0,a]$, the fundamental theorem gives
\[
 |v(x',0)|^p
 \leq |v(x',t)|^p+
 p\int_0^a|v|^{p-1}|\partial_n v|(x',s)\,ds.
\]
For $p=1$, use the absolutely continuous function $|v|$ and its
derivative bound $|(|v|)'|\leq|v'|$.
Average $t$ over $(0,a)$ and integrate $x'$.
For $p>1$, Young's inequality bounds
$p|v|^{p-1}|\partial_n v|$ by $C_p(|v|^p+|Dv|^p)$.
Consequently
\[
 \int_{Q'}|v(x',0)|^p\,dx'
 \leq C_{p,a}\int_{Q'\times(0,a)}(|v|^p+|Dv|^p)\,dx.
\]
Boundary flattening introduces only bounded Jacobian and derivative
factors, including $\sqrt{1+|Dg|^2}$ for surface measure.
A finite partition yields the global estimate for smooth $u$.

\paragraph{2. Complete the restriction map.}
Choose $u_j$ smooth near $\overline U$, with $u_j\to u$ in $W^{1,p}$.
The estimate on $u_j-u_l$ shows that $u_j|_{\partial U}$ is Cauchy in
$L^p(\partial U)$. Define $Tu$ to be its limit. The same estimate on
two approximating sequences proves independence of choice and boundedness.
Linearity and uniqueness follow from density.

\paragraph{3. Agreement with continuous functions.}
For $u$ continuous on $\overline U$, the inward-translation and smoothing
construction in \S\ref{sec:5-3} can be chosen to converge uniformly as
well as in $W^{1,p}$. Uniform continuity on the compact closure justifies
this assertion for each localized piece. The boundary restrictions
therefore converge to $u|_{\partial U}$ uniformly, identifying the trace.
\end{expandedproof}
\Needspace{45mm}
\PDEAnnotation{迹估计的实际用处}{只要 $u_j\to u$ 于 $W^{1,p}$，就自动有 $Tu_j\to Tu$ 于边界 $L^p$。边界条件因此可随强逼近传递，而不要求逐点追踪边界值。}

\subsection{Zero trace and the closure definition}
\begin{theorem}[Evans \S5.5, Theorem 2]\label{thm:5-zero-trace}
Under the same hypotheses,
\[
 W^{1,p}_0(U)=\{u\in W^{1,p}(U):Tu=0\}.
\]
\end{theorem}
\begin{expandedproof}
If $u_j\in C_c^\infty(U)$ tends to $u$ in $W^{1,p}$, boundedness of $T$
gives $Tu=\lim Tu_j=0$.

For the reverse direction we supply the boundary-layer argument.
In a flat chart the trace equals the value at $t=0$ of the absolutely
continuous representative on almost every vertical line. To justify
this, use smooth approximations, Fubini, and the one-dimensional estimate
\[
 |w(0)|\leq a^{-1}\|w\|_{L^1(0,a)}+\|w'\|_{L^1(0,a)}.
\]
A subsequence converges in $W^{1,p}$ on almost every line, and this
estimate identifies the endpoint limit with the trace limit.

Thus if $Tv=0$,
$v(x',t)=\int_0^t\partial_n v(x',s)\,ds$ on almost every such line.
H\"older (or the integral triangle inequality for $p=1$) gives
\[
 \int_{Q'}\int_0^{2\delta}|v(x',t)|^p\,dt\,dx'
 \leq (2\delta)^p
       \int_{Q'}\int_0^{2\delta}|\partial_n v(x',t)|^p\,dt\,dx'.
\]
Choose a smooth $\chi_\delta(t)$ equal to zero for $t<\delta$, one
for $t>2\delta$, with $|\chi_\delta'|\leq C/\delta$. Then
\[
 D(\chi_\delta v)-Dv
 =(\chi_\delta-1)Dv+vD\chi_\delta.
\]
The first term tends to zero by absolute continuity of the integral.
The $p$-th power of the second norm is bounded by
$C\int_{0<t<2\delta}|Dv|^p$, hence also tends to zero.
The function error tends to zero in the same way.
Flattening and a finite partition transfer this argument to $U$.
It gives functions supported a positive distance from $\partial U$
converging to $u$ in $W^{1,p}$. For each one, convolution with a smaller
radius produces an element of $C_c^\infty(U)$ with arbitrarily small
additional error. A diagonal choice completes the proof.
\end{expandedproof}
\Needspace{45mm}
\PDEAnnotation{截断产生的危险项}{$D(\chi_\delta u)$ 中有 $uD\chi_\delta$，而 $D\chi_\delta$ 约为 $1/\delta$。零迹给出的法向积分公式补回一个 $\delta$ 因子，才使此项趋于零。}

\paragraph{Zero extension.}
If $u\in W^{1,p}_0(U)$, extend compactly supported approximations by
zero. They are Cauchy in $W^{1,p}(\R^n)$ with exactly the same norms,
so completeness gives a zero extension $\widetilde u$ satisfying
$D\widetilde u=\widetilde{Du}$. This argument works without regularity
of $U$ for finite $p$. A nonzero trace generally produces a boundary
distribution instead.
